% GENERATED FILE. Edit canonical lessons, then run npm run generate:books.
% canonical-source-hash: fnv1a64:753a761452bf5dc8
% canonical-lessons: TE-C464-andincu, TE-C464-samakurcu, TE-C464-erpatu-ceyu, TE-C464-saraphara, TE-C464-vasulu-ceyu

\chapter{Arranging Things}
\label{ch:te-arranging-things}
\hlchaptermodality{464}

\begin{chapteropening}
\textbf{By the end of this chapter:} \emph{I can say to provide, to gather together, to organise, a supply and to collect payment.}

Complete the last lesson of chapter 464: I can say to provide, to gather together, to organise, a supply and to collect payment.
\end{chapteropening}

\section[\texorpdfstring{\te{అందించు} (andiñcu)}{అందించు (andiñcu)}]{\te{అందించు} (andiñcu) — to provide; to hand over}

\label{lesson:TE-C464-andincu}



\begin{quote}
Before the new one: say the Telugu for to extend, then the Telugu for to renew.
\end{quote}

\subsection*{You'll want to know: \te{అందించు}}
\textbf{\te{అందించు}} — \emph{andiñcu} — \textquotedblleft{}to provide; to hand over\textquotedblright{}.

\begin{quote}
\te{ఆ} \te{సత్రం} \te{ప్రయాణికులకు} \te{ఉచితంగా} \te{భోజనం} \te{అందించేది}.
\end{quote}

\begin{quote}
\emph{ā satraṁ prayāṇikulaku ucitaṅgā bhōjanaṁ andiñcēdi.}
\end{quote}

\begin{quote}
That inn used to provide travellers with free meals.
\end{quote}

Past (I …): \te{అందించాను}, \emph{andiñcānu}.

The causative of \te{అందు}, to reach: to make something reach someone. The receiver takes -\te{కి}.

\begin{cousinweb}
\textbf{\te{అందించు}} (\emph{andiñcu}) is native Dravidian: \te{అందు}, \textquotedblleft{}to reach, to be within reach\textquotedblright{} (DEDR 149), made causative with -\te{ఇంచు}.
\end{cousinweb}

\begin{grammarlens}[title={What you've built}]
One more word for getting around.
\end{grammarlens}

\subsection*{Your turn}
\begin{itemize}
\raggedright
  \item \emph{Say it:} \emph{andiñcu}
  \item \emph{Say it:} \emph{andiñcu}, once more
  \item \emph{Say it:} \emph{ā satraṁ prayāṇikulaku ucitaṅgā bhōjanaṁ andiñcēdi.} — That inn used to provide travellers with free meals.
  \item \emph{Recall:} say the Telugu for to extend, then the Telugu for to renew, then say \emph{andiñcu} again
\end{itemize}

\begin{tcolorbox}[breakable,colback=teal!4,colframe=teal!35!black,title={Before you move on}]
What does \te{అందించు} mean? (\textquotedblleft{}To provide; to hand over\textquotedblright{}.) Say it once more.
\end{tcolorbox}
\section[\texorpdfstring{\te{సమకూర్చు} (samakūrcu)}{సమకూర్చు (samakūrcu)}]{\te{సమకూర్చు} (samakūrcu) — to gather together; to arrange}

\label{lesson:TE-C464-samakurcu}



\begin{quote}
Before the new one: say the Telugu for to renew, then the Telugu for to provide.
\end{quote}

\subsection*{You'll want to know: \te{సమకూర్చు}}
\textbf{\te{సమకూర్చు}} — \emph{samakūrcu} — \textquotedblleft{}to gather together; to arrange\textquotedblright{}.

\begin{quote}
\te{పెళ్ళి} \te{కోసం} \te{నాన్న} \te{డబ్బు} \te{సమకూర్చాడు}.
\end{quote}

\begin{quote}
\emph{peḷḷi kōsaṁ nānna ḍabbu samakūrcāḍu.}
\end{quote}

\begin{quote}
Father got the money together for the wedding.
\end{quote}

Past (I …): \te{సమకూర్చాను}, \emph{samakūrcānu}.

Said of money, things or people brought together for a purpose.

\begin{cousinweb}
\textbf{\te{సమకూర్చు}} (\emph{samakūrcu}) is a Telugu verb: \te{కూర్చు}, \textquotedblleft{}to join, to bring together, to collect\textquotedblright{}, native Dravidian (DEDR 1882, beside \te{కూడు}, \textquotedblleft{}to join\textquotedblright{}), with \te{సమ}-, usually read as Sanskrit \emph{sama}, \textquotedblleft{}together, even\textquotedblright{}.
\end{cousinweb}

\begin{grammarlens}[title={What you've built}]
One more word for getting around.
\end{grammarlens}

\subsection*{Your turn}
\begin{itemize}
\raggedright
  \item \emph{Say it:} \emph{samakūrcu}
  \item \emph{Say it:} \emph{samakūrcu}, once more
  \item \emph{Say it:} \emph{peḷḷi kōsaṁ nānna ḍabbu samakūrcāḍu.} — Father got the money together for the wedding.
  \item \emph{Recall:} say the Telugu for to renew, then the Telugu for to provide, then say \emph{samakūrcu} again
\end{itemize}

\begin{tcolorbox}[breakable,colback=teal!4,colframe=teal!35!black,title={Before you move on}]
What does \te{సమకూర్చు} mean? (\textquotedblleft{}To gather together; to arrange\textquotedblright{}.) Say it once more.
\end{tcolorbox}
\section[\texorpdfstring{\te{ఏర్పాటు} \te{చేయు} (ērpāṭu cēyu)}{ఏర్పాటు చేయు (ērpāṭu cēyu)}]{\te{ఏర్పాటు} \te{చేయు} (ērpāṭu cēyu) — to organise; to set up}

\label{lesson:TE-C464-erpatu-ceyu}



\begin{quote}
Before the new one: say the Telugu for to provide, then the Telugu for to gather together.
\end{quote}

\subsection*{You'll want to know: \te{ఏర్పాటు} \te{చేయు}}
\textbf{\te{ఏర్పాటు} \te{చేయు}} — \emph{ērpāṭu cēyu} — \textquotedblleft{}to organise; to set up\textquotedblright{}.

\begin{quote}
\te{గ్రామంలో} \te{వైద్య} \te{శిబిరం} \te{ఏర్పాటు} \te{చేశారు}.
\end{quote}

\begin{quote}
\emph{grāmaṁlō vaidya śibiraṁ ērpāṭu cēśāru.}
\end{quote}

\begin{quote}
They set up a medical camp in the village.
\end{quote}

Past (I …): \te{ఏర్పాటు} \te{చేశాను}, \emph{ērpāṭu cēśānu}.

\te{ఏర్పాటు}, an arrangement, + \te{చేయు}. \te{ఏర్పాట్లు}, the arrangements for an event.

\begin{cousinweb}
\textbf{\te{ఏర్పాటు} \te{చేయు}} (\emph{ērpāṭu cēyu}) is native Dravidian: \te{ఏర్పాటు}, \textquotedblleft{}an arrangement\textquotedblright{}, the noun of \te{ఏర్పడు}, \textquotedblleft{}to be arranged, to be established\textquotedblright{} (DEDR 905), with \te{చేయు}, \textquotedblleft{}to do\textquotedblright{} (DEDR 1957).
\end{cousinweb}

\begin{grammarlens}[title={What you've built}]
One more word for getting around.
\end{grammarlens}

\subsection*{Your turn}
\begin{itemize}
\raggedright
  \item \emph{Say it:} \emph{ērpāṭu cēyu}
  \item \emph{Say it:} \emph{ērpāṭu cēyu}, once more
  \item \emph{Say it:} \emph{grāmaṁlō vaidya śibiraṁ ērpāṭu cēśāru.} — They set up a medical camp in the village.
  \item \emph{Recall:} say the Telugu for to provide, then the Telugu for to gather together, then say \emph{ērpāṭu cēyu} again
\end{itemize}

\begin{tcolorbox}[breakable,colback=teal!4,colframe=teal!35!black,title={Before you move on}]
What does \te{ఏర్పాటు} \te{చేయు} mean? (\textquotedblleft{}To organise; to set up\textquotedblright{}.) Say it once more.
\end{tcolorbox}
\section[\texorpdfstring{\te{సరఫరా} (sarapharā)}{సరఫరా (sarapharā)}]{\te{సరఫరా} (sarapharā) — a supply}

\label{lesson:TE-C464-saraphara}



\begin{quote}
Before the new one: say the Telugu for to gather together, then the Telugu for to organise.
\end{quote}

\subsection*{You'll want to know: \te{సరఫరా}}
\textbf{\te{సరఫరా}} — \emph{sarapharā} — \textquotedblleft{}a supply\textquotedblright{}.

\begin{quote}
\te{రెండు} \te{రోజులుగా} \te{నీటి} \te{సరఫరా} \te{లేదు}.
\end{quote}

\begin{quote}
\emph{reṇḍu rōjulugā nīṭi sarapharā lēdu.}
\end{quote}

\begin{quote}
There has been no water supply for two days.
\end{quote}

\te{నీటి} \te{సరఫరా}, water supply; \te{విద్యుత్} \te{సరఫరా}, power supply. \te{సరఫరా} \te{చేయు}, to supply.

\begin{cousinweb}
\textbf{\te{సరఫరా}} (\emph{sarapharā}) is usually traced to Persian \emph{sarbarāh}, \textquotedblleft{}a manager; provision\textquotedblright{}, which reached Telugu through Hindustani.
\end{cousinweb}

\begin{grammarlens}[title={What you've built}]
One more word for getting around.
\end{grammarlens}

\subsection*{Your turn}
\begin{itemize}
\raggedright
  \item \emph{Say it:} \emph{sarapharā}
  \item \emph{Say it:} \emph{sarapharā}, once more
  \item \emph{Say it:} \emph{reṇḍu rōjulugā nīṭi sarapharā lēdu.} — There has been no water supply for two days.
  \item \emph{Recall:} say the Telugu for to gather together, then the Telugu for to organise, then say \emph{sarapharā} again
\end{itemize}

\begin{tcolorbox}[breakable,colback=teal!4,colframe=teal!35!black,title={Before you move on}]
What does \te{సరఫరా} mean? (\textquotedblleft{}A supply\textquotedblright{}.) Say it once more.
\end{tcolorbox}
\section[\texorpdfstring{\te{వసూలు} \te{చేయు} (vasūlu cēyu)}{వసూలు చేయు (vasūlu cēyu)}]{\te{వసూలు} \te{చేయు} (vasūlu cēyu) — to collect payment; to charge}

\label{lesson:TE-C464-vasulu-ceyu}



\begin{quote}
Before the new one: say the Telugu for to organise, then the Telugu for a supply.
\end{quote}

\subsection*{You'll want to know: \te{వసూలు} \te{చేయు}}
\textbf{\te{వసూలు} \te{చేయు}} — \emph{vasūlu cēyu} — \textquotedblleft{}to collect payment; to charge\textquotedblright{}.

\begin{quote}
\te{ఆటోవాడు} \te{మా} \te{దగ్గర} \te{రెట్టింపు} \te{కిరాయి} \te{వసూలు} \te{చేశాడు}.
\end{quote}

\begin{quote}
\emph{āṭōvāḍu mā daggara reṭṭimpu kirāyi vasūlu cēśāḍu.}
\end{quote}

\begin{quote}
The auto driver charged us double the fare.
\end{quote}

Past (I …): \te{వసూలు} \te{చేశాను}, \emph{vasūlu cēśānu}.

Said of fares, fees, rent and taxes collected from people.

\begin{cousinweb}
\textbf{\te{వసూలు} \te{చేయు}} (\emph{vasūlu cēyu}) is a Telugu phrase: \te{వసూలు}, usually traced to Arabic \emph{wuṣūl}, \textquotedblleft{}arrival, receipt\textquotedblright{}, through Persian and Hindustani, and the native \te{చేయు}, \textquotedblleft{}to do\textquotedblright{} (DEDR 1957).
\end{cousinweb}

\begin{grammarlens}[title={What you've built}]
One more word for getting around.
\end{grammarlens}

\subsection*{Your turn}
\begin{itemize}
\raggedright
  \item \emph{Say it:} \emph{vasūlu cēyu}
  \item \emph{Say it:} \emph{vasūlu cēyu}, once more
  \item \emph{Say it:} \emph{āṭōvāḍu mā daggara reṭṭimpu kirāyi vasūlu cēśāḍu.} — The auto driver charged us double the fare.
  \item \emph{Recall:} say the Telugu for to organise, then the Telugu for a supply, then say \emph{vasūlu cēyu} again
\end{itemize}

\begin{tcolorbox}[breakable,colback=teal!4,colframe=teal!35!black,title={Before you move on}]
What does \te{వసూలు} \te{చేయు} mean? (\textquotedblleft{}To collect payment; to charge\textquotedblright{}.) Say it once more.
\end{tcolorbox}
